% Provenance (never printed): sources/book2/BRIEFS.md, g12-11. Budget and primary balance; automatic stabilisers; the
% debt-ratio recurrence and its limit; sustainability; Ricardian equivalence (first look); how large debts ended; fiscal rules
% in a currency union. Real numbers: figdata/grade-12/11-fiscal-policy-and-debt/ (JST debtgdp, gdp, ltrate, cpi for GBR and
% USA; Eurostat gov_10dd_edpt1 GD and B9). Ledger: TFEU Protocol 12 (3 % and 60 %), Parker et al. 2013 (abstract of the working
% paper), Wikipedia on the Greek debt crisis (the 53.5 % face-value cut; second programme ratified February 2012). The
% Reinhart-Rogoff table of how debts ended was not read: the chapter tells the exits through its own two records.
\chapter{Fiscal Policy and Public Debt}\label{ch:g12:fiscal-policy-and-debt}

In 1946 the British government owed 270\,\% of a year's national income. Thirty years later it owed 47\,\%. It did not repay most of its debt: in money, the debt grew.
What fell was the debt as a share of a national income that grew much faster.
% figdata: debt_key (uk1946 269.8, uk1975 46.7; nominal debt 1975 / 1946 = uk_nominal_ratio)

This chapter is about a government's budget over time: how deficits add to debt, what makes a \omterm{def:g12:fiscal-policy-and-debt:debt-ratio}{debt ratio} rise or fall, when a debt can be carried, and how the largest
debts of the past came down.

\begin{openingchart}{Public debt as a percentage of GDP, United Kingdom and United States, 1870--2020.}{Jordà--Schularick--Taylor Macrohistory Database}
\begin{tikzpicture}
\begin{axis}[omtime, width=0.9\linewidth, height=6.0cm, xmin=1870, xmax=2020, ymin=0, ymax=280, ytick={0,50,...,250}, ylabel={per cent of GDP},
  xtick={1870,1890,...,2010}]
\addplot[omDef, thick] table[x=year, y=gbr] {figdata/grade-12/11-fiscal-policy-and-debt/out/debt_long.csv};
\addplot[omThm, thick] table[x=year, y=usa] {figdata/grade-12/11-fiscal-policy-and-debt/out/debt_long.csv};
\node[font=\scriptsize, omDef, anchor=west] at (axis cs:1947,262) {United Kingdom};
\node[font=\scriptsize, omThm, anchor=west] at (axis cs:1890,22) {United States};
\end{axis}
\end{tikzpicture}
\end{openingchart}

\begin{recall}
From the grades 6--9 volume: the government's budget, deficits and debt, and the ageing of populations. From \cref{ch:g12:present-value}: \omterm{def:g12:present-value:present-value}{present value}. From
\cref{ch:g12:saving-and-investment}: \omterm{def:g12:saving-and-investment:crowding-out}{crowding out}. From \cref{ch:g12:keynesian-cross}: the multiplier. From mathematics: sequences of the form $u_{t+1} = a u_t + b$.
\end{recall}

\section{The budget balance}\label{sec:g12:fiscal-policy-and-debt:balance}

\begin{definition}[Primary balance]\label{def:g12:fiscal-policy-and-debt:primary-balance}
The \emph{primary balance}\index{primary balance} is government revenue minus government spending other than interest on the debt.
\end{definition}

The overall balance is the \omterm{def:g12:fiscal-policy-and-debt:primary-balance}{primary balance} minus interest. A government can run a primary surplus and an overall deficit at once, if its interest bill is large.

\begin{definition}[Automatic stabiliser]\label{def:g12:fiscal-policy-and-debt:automatic-stabiliser}
An \emph{automatic stabiliser}\index{automatic stabiliser} is a part of the budget that, without any new decision, takes in less or pays out more when output falls, such as
taxes on income and profits or benefits for the unemployed.
\end{definition}

\begin{definition}[Structural balance]\label{def:g12:fiscal-policy-and-debt:structural-balance}
The \emph{structural balance}\index{structural balance} is the budget balance the government would have if output were at potential: the actual balance with the effect of
the cycle removed. Like \omterm{def:g12:aggregate-demand-and-supply:potential-output}{potential output}, it is estimated, not observed.
\end{definition}

In the euro area as a whole, the general government deficit went from 0.7\,\% of GDP in 2007 to 6.3\,\% in 2009, as tax receipts fell with output and spending on
benefits and rescues rose.
% figdata: debt_key (ea_b9_2007 -0.7, ea_b9_2009 -6.3)

\section{How the debt ratio moves}\label{sec:g12:fiscal-policy-and-debt:dynamics}

\begin{definition}[Debt ratio]\label{def:g12:fiscal-policy-and-debt:debt-ratio}
The \emph{debt ratio}\index{debt ratio} is the government's debt divided by a year's GDP.
\end{definition}

\begin{proposition}[The debt recurrence]\label{prop:g12:fiscal-policy-and-debt:recurrence}
Let $d_t$ be the \omterm{def:g12:fiscal-policy-and-debt:debt-ratio}{debt ratio}, $i$ the interest rate paid on the debt, $\gamma$ the growth rate of \omterm{def:g10:national-accounts:nominal-gdp}{nominal GDP} and $p$ the primary surplus as a share of GDP. Then
\[ d_{t+1} = \frac{1 + i}{1 + \gamma}\, d_t - p. \]
\end{proposition}
\begin{proof}
Debt next year is debt this year plus interest minus the primary surplus: $D_{t+1} = (1 + i) D_t - P_{t+1}$. Divide by $Y_{t+1} = (1 + \gamma) Y_t$:
$d_{t+1} = (1 + i) D_t / ((1 + \gamma) Y_t) - P_{t+1}/Y_{t+1} = \frac{1 + i}{1 + \gamma} d_t - p$.
\end{proof}

Write $a = (1 + i)/(1 + \gamma)$. When the interest rate is below the growth rate, $a < 1$: the ratio shrinks by itself unless the government runs primary deficits. When
it is above, $a > 1$: the ratio grows unless the government runs primary surpluses. The \omterm{def:g12:fiscal-policy-and-debt:primary-balance}{primary balance} that holds the ratio constant is $p = (a - 1) d = d\,(i -
\gamma)/(1 + \gamma)$.

\begin{omfigure}
\begin{tikzpicture}
\begin{groupplot}[group style={group size=2 by 1, horizontal sep=1.2cm}, width=0.47\linewidth, height=5.2cm, xmin=0, xmax=110, ymin=0, ymax=110,
  xlabel={$d_t$}, tick label style={font=\scriptsize}, label style={font=\scriptsize}, xtick={0,40,100}, ytick={0,40,100}]
\nextgroupplot[title={\scriptsize $a = 0.8$, primary deficit 8: converges}, ylabel={$d_{t+1}$}]
\addplot[gray, domain=0:110, samples=2] {x};
\addplot[omDef, thick, domain=0:110, samples=2] {0.8*x + 8};
\addplot[omThm, thin] coordinates {(100,0) (100,88) (88,88) (88,78.4) (78.4,78.4) (78.4,70.72) (70.72,70.72) (70.72,64.58) (64.58,64.58) (64.58,59.66)};
\fill (axis cs:40,40) circle (1.5pt);
\nextgroupplot[title={\scriptsize $a = 1.2$, primary surplus 8: diverges}]
\addplot[gray, domain=0:110, samples=2] {x};
\addplot[omDef, thick, domain=0:110, samples=2] {1.2*x - 8};
\addplot[omThm, thin] coordinates {(50,0) (50,52) (52,52) (52,54.4) (54.4,54.4) (54.4,57.28) (57.28,57.28) (57.28,60.74) (60.74,60.74) (60.74,64.88)};
\fill (axis cs:40,40) circle (1.5pt);
\end{groupplot}
\end{tikzpicture}

{\small The recurrence drawn step by step, with numbers exaggerated for the picture. Left: from a ratio of 100, each year moves up to the blue line and across to the grey
line $d_{t+1} = d_t$; the ratio falls towards 40. Right: the resting point is again 40, but a ratio that starts at 50 moves away from it, faster each year.}
\end{omfigure}

\begin{definition}[Sustainable debt]\label{def:g12:fiscal-policy-and-debt:sustainable}
A government's debt is \emph{sustainable}\index{sustainable debt} when the \omterm{def:g12:fiscal-policy-and-debt:primary-balance}{primary balances} it can be expected to run keep the \omterm{def:g12:fiscal-policy-and-debt:debt-ratio}{debt ratio} from rising without limit, without
default or \omterm{def:g12:inflation-and-money:unexpected-inflation}{unexpected inflation}.
\end{definition}

\begin{model}{Debt dynamics}
\assumes The recurrence (an identity) with given interest rate, growth rate and \omterm{def:g12:fiscal-policy-and-debt:primary-balance}{primary balance}.
\predicts The path of the \omterm{def:g12:fiscal-policy-and-debt:debt-ratio}{debt ratio}; a stable ratio when $i < \gamma$, an explosive one when $i > \gamma$ unless primary surpluses offset it.
\tested Consistent with it, as an accounting relation must be: the United Kingdom after 1945 and the United States after 1946, ratios above 100\,\% fell for thirty years with nominal growth above the interest rate (the opening chart
  and the weekend problem).
\breaks The interest rate depends on the debt: when lenders fear default, they ask for a higher rate, which raises $a$ and the debt it was meant to measure, as in Greece in
  2010--12. Growth too may fall when the primary surplus is raised to stabilise the debt, through the multiplier.
\end{model}

\section{Do households look through the debt?}\label{sec:g12:fiscal-policy-and-debt:ricardian}

\begin{definition}[Ricardian equivalence]\label{def:g12:fiscal-policy-and-debt:ricardian-equivalence}
\emph{Ricardian equivalence}\index{Ricardian equivalence} is the proposition that, for given government spending, replacing taxes today by borrowing does not change
households' consumption, because they save the tax cut to pay the future taxes that will repay the debt.
\end{definition}

\begin{model}{Ricardian equivalence}
\assumes Households look ahead to future taxes, can borrow and lend freely, and care about the taxes their heirs will pay.
\predicts A debt-financed tax cut or rebate is saved, not spent.
\tested In 2008 American households received one-off stimulus payments. In the three months of receipt they spent about 12--30\,\% of the payments on nondurable goods, and about 50--90\,\% counting durables such as cars. Much of the money was spent, not saved.
\breaks Households that cannot borrow, or do not look far ahead, spend a rebate; the proposition holds best for those who were already saving.
\end{model}
% ledger: g12:parker2013

\section{How large debts ended}\label{sec:g12:fiscal-policy-and-debt:exits}

The recurrence shows the four ways a \omterm{def:g12:fiscal-policy-and-debt:debt-ratio}{debt ratio} can come down: growth of real output, inflation, primary surpluses, and default, which cuts $D$ itself. A government that
borrows in its own currency can always pay in money its central bank creates; the record of the hyperinflations in the grades 6--9 volume shows the price of that exit.

\begin{record}[How the United Kingdom's debt fell after 1945]
% figdata: uk_decomp (interest 189.5; real growth -104.8; inflation -192.2; other -115.7); debt_key (gamma_mean 9, i_mean 6.42)
From 1946 to 1975 the British \omterm{def:g12:fiscal-policy-and-debt:debt-ratio}{debt ratio} fell from 269.8 to 46.7\,\% of GDP. \omterm{def:g10:national-accounts:nominal-gdp}{Nominal GDP} grew by 9.0\,\% a year on average and the long-term interest rate averaged 6.4\,\%.
Year by year, the recurrence splits the fall of 223 points into interest, which added 190 points; real growth, which removed 105; inflation, which removed 192; and the
rest, primary surpluses and the gap between the long-term rate and the rate actually paid on the debt, which removed 116. The weekend problem works through it.
\end{record}

\begin{definition}[Sovereign default]\label{def:g12:fiscal-policy-and-debt:default}
A \emph{sovereign default}\index{sovereign default} is a government's failure to pay interest or principal on its debt as promised, including a forced exchange of its \omterm{def:g12:present-value:bond}{bonds}
for \omterm{def:g12:present-value:bond}{bonds} worth less.
\end{definition}

\begin{record}[Greece 2010--12]
% ledger: g12:greece-psi; figdata: debt_key (el_gd2008 110.9, el_b9_2009 -15.4, el_gd2011 175.1, el_gd2012 164.1)
Greece's general government deficit reached 15.4\,\% of GDP in 2009, and its debt rose from 110.9\,\% of GDP in 2008 to 175.1\,\% in 2011. Lenders demanded higher rates, and the
government borrowed from other euro-area governments and the IMF instead. The second rescue programme, ratified in February 2012, required private bondholders to accept a
53.5\,\% reduction in the face value of their \omterm{def:g12:present-value:bond}{bonds}. The ratio fell to 164.1\,\% in 2012, then rose again, to 180.4\,\% in 2013.
\end{record}

\begin{definition}[Fiscal rule]\label{def:g12:fiscal-policy-and-debt:fiscal-rule}
A \emph{fiscal rule}\index{fiscal rule} is a lasting limit on a budget figure, such as the deficit, the debt or spending, set in law.
\end{definition}

\begin{omchart}{General government gross debt, per cent of GDP, six euro-area countries, 2000--2025.}{Eurostat (gov\_10dd\_edpt1)}
\begin{tikzpicture}
\begin{axis}[omtime, width=0.9\linewidth, height=5.6cm, xmin=2000, xmax=2025, ymin=0, ymax=220, ytick={0,60,100,150,200}, ylabel={per cent of GDP}]
\addplot[gray, dashed, domain=2000:2025, samples=2] {60};
\addplot[omDef, thick] table[x=year, y=EL] {figdata/grade-12/11-fiscal-policy-and-debt/out/debt_euro.csv};
\addplot[omThm, thick] table[x=year, y=IT] {figdata/grade-12/11-fiscal-policy-and-debt/out/debt_euro.csv};
\addplot[omProp, thick] table[x=year, y=FR] {figdata/grade-12/11-fiscal-policy-and-debt/out/debt_euro.csv};
\addplot[omMeth, thick] table[x=year, y=ES] {figdata/grade-12/11-fiscal-policy-and-debt/out/debt_euro.csv};
\addplot[black, thick] table[x=year, y=DE] {figdata/grade-12/11-fiscal-policy-and-debt/out/debt_euro.csv};
\addplot[gray, thick] table[x=year, y=IE] {figdata/grade-12/11-fiscal-policy-and-debt/out/debt_euro.csv};
\node[font=\scriptsize, omDef, anchor=south] at (axis cs:2018,186) {Greece};
\node[font=\scriptsize, omThm, anchor=south] at (axis cs:2005,108) {Italy};
\node[font=\scriptsize, omProp, anchor=west] at (axis cs:2022.3,118) {France};
\node[font=\scriptsize, omMeth, anchor=north] at (axis cs:2004,40) {Spain};
\node[font=\scriptsize, anchor=south] at (axis cs:2017,66) {Germany};
\node[font=\scriptsize, gray, anchor=north] at (axis cs:2022,34) {Ireland};
\end{axis}
\end{tikzpicture}
\end{omchart}

\begin{record}[A fiscal rule in a currency union]
% ledger: g12:tfeu-p12; figdata: debt_key (euro12_over3 124, euro12_obs 288, gd2025_*)
The treaties of the European Union set reference values of 3\,\% of GDP for the government deficit and 60\,\% for the government debt. Among the twelve countries that have
used the euro since 2002 or earlier, the deficit exceeded 3\,\% of GDP in 124 of the 288 country-years from 2002 to 2025, 43\,\%. In 2025 the debt was below 60\,\% in
Ireland (34.8\,\%), close to it in Germany (62.8\,\%), and far above it in Spain (100.5\,\%), France (115.7\,\%), Italy (136.7\,\%) and Greece (145.4\,\%).
\end{record}

\section{Exercises}

\begin{exercise}[$\star$]\label{exo:g12:fiscal-policy-and-debt:1}
A government's deficit is 4\,\% of GDP, and its interest payments are 3\,\% of GDP. Find its \omterm{def:g12:fiscal-policy-and-debt:primary-balance}{primary balance}.
\end{exercise}

\begin{exercise}[$\star$]\label{exo:g12:fiscal-policy-and-debt:2}
Revenue is 40, spending other than interest 38 and interest 5, all in per cent of GDP. Find the \omterm{def:g12:fiscal-policy-and-debt:primary-balance}{primary balance} and the overall balance.
\end{exercise}

\begin{exercise}[$\star$]\label{exo:g12:fiscal-policy-and-debt:3}
From the opening chart, when did the British \omterm{def:g12:fiscal-policy-and-debt:debt-ratio}{debt ratio} peak, and at what level? What was it in 1975?
\end{exercise}

\begin{exercise}[$\star$]\label{exo:g12:fiscal-policy-and-debt:4}
From the opening chart, compare the American \omterm{def:g12:fiscal-policy-and-debt:debt-ratio}{debt ratio} in 1946 and in 2020.
\end{exercise}

\begin{exercise}[$\star$]\label{exo:g12:fiscal-policy-and-debt:5}
The \omterm{def:g12:fiscal-policy-and-debt:debt-ratio}{debt ratio} is 60\,\%, the interest rate 4\,\% and nominal growth 3\,\%, with a balanced primary budget. Find next year's ratio.
\end{exercise}

\begin{exercise}[$\star\star$]\label{exo:g12:fiscal-policy-and-debt:6}
From a ratio of 100\,\%, with $i = 2\,\%$, $\gamma = 5\,\%$ and a primary deficit of 1\,\% of GDP, find the ratio after one, two and three years.
\end{exercise}

\begin{exercise}[$\star\star$]\label{exo:g12:fiscal-policy-and-debt:7}
From a ratio of 100\,\%, with $i = 6\,\%$, $\gamma = 3\,\%$ and a balanced primary budget, find the ratio after one, two and three years.
\end{exercise}

\begin{exercise}[$\star\star$]\label{exo:g12:fiscal-policy-and-debt:8}
The ratio is 120\,\%, $i = 5\,\%$ and $\gamma = 3\,\%$. Find the \omterm{def:g12:fiscal-policy-and-debt:primary-balance}{primary balance} that keeps the ratio constant.
\end{exercise}

\begin{exercise}[$\star\star$]\label{exo:g12:fiscal-policy-and-debt:9}
The ratio is 100\,\% and $\gamma = 4\,\%$. The interest rate rises from 3\,\% to 5\,\%. How does the stabilising \omterm{def:g12:fiscal-policy-and-debt:primary-balance}{primary balance} change?
\end{exercise}

\begin{exercise}[$\star\star$]\label{exo:g12:fiscal-policy-and-debt:10}
From the chart of euro-area debt, which country's ratio rose most from 2007 to 2025, and which fell most after 2012? Use the record's values for 2025 and 2007 values of
35.7\,\% (Spain), 65.5\,\% (France) and 24.0\,\% (Ireland).
\end{exercise}

\begin{exercise}[$\star\star$]\label{exo:g12:fiscal-policy-and-debt:11}
Tax receipts are 40\,\% of GDP and move in proportion to it, and benefits rise by 0.3\,\% of GDP when output falls by 3\,\%. By how much does the deficit widen, in per cent
of GDP, with no new decision?
\end{exercise}

\begin{exercise}[$\star\star\star$]\label{exo:g12:fiscal-policy-and-debt:12}
What evidence contradicts \omterm{def:g12:fiscal-policy-and-debt:ricardian-equivalence}{Ricardian equivalence}? Which assumption of the model does it point to?
\end{exercise}

\begin{exercise}[$\star\star\star$]\label{exo:g12:fiscal-policy-and-debt:13}
Let $d^*$ satisfy $d^* = a d^* - p$, with $a \neq 1$. Show that $d_t - d^* = a^t (d_0 - d^*)$, and deduce when the ratio converges.
\end{exercise}

\begin{exercise}[$\star\star\star$]\label{exo:g12:fiscal-policy-and-debt:14}
A government owes 100\,\% of GDP in long-term \omterm{def:g12:present-value:bond}{bonds} at a fixed rate. Prices rise by 10\,\% more than expected in one year, with real output unchanged. Ignoring interest and
the \omterm{def:g12:fiscal-policy-and-debt:primary-balance}{primary balance}, what is the new ratio? Who loses?
\end{exercise}

\begin{exercise}[$\star\star\star$]\label{exo:g12:fiscal-policy-and-debt:15}
A government owes 170\,\% of GDP, half of it to private bondholders, who accept a cut of 53.5\,\% in the face value of their \omterm{def:g12:present-value:bond}{bonds}. Find the new ratio. Why might it rise again
afterwards?
\end{exercise}

\section{Problem: How the United Kingdom's Debt Ratio Fell, 1946--1975}

\begin{problem}[{Weekend problem --- growth, inflation and surpluses}]
\label{pb:g12:fiscal-policy-and-debt:1}
From 1946 to 1975 the British \omterm{def:g12:fiscal-policy-and-debt:debt-ratio}{debt ratio} fell from 269.8\,\% to 46.7\,\% of GDP. Over 1947--1975, \omterm{def:g10:national-accounts:nominal-gdp}{nominal GDP} grew by 9.0\,\% a year on average, the long-term interest rate
averaged 6.42\,\%, consumer price inflation 5.82\,\% and real growth 3.05\,\%. The table splits the change in the ratio, year by year with the recurrence, using the long-term
rate as the interest rate on the debt (in percentage points of GDP, summed over the years; share of the fall in per cent, computed from unrounded
values).

\medskip
\begin{center}
\small
\pgfplotstabletypeset[col sep=comma, columns={item,points,share},
  columns/item/.style={string type, column name={contribution}, column type=l},
  columns/points/.style={column name={points}, fixed, precision=1, zerofill},
  columns/share/.style={column name={share of the fall}, fixed, precision=1, zerofill},
  every head row/.style={before row=\toprule, after row=\midrule}, every last row/.style={after row=\bottomrule}]
  {figdata/grade-12/11-fiscal-policy-and-debt/out/uk_decomp.csv}
\end{center}

\medskip
\noindent\textbf{Part I --- The averages.}
\begin{enumerate}
\item By how many points did the ratio fall?
\item Check that real growth and inflation roughly account for nominal growth.
\item With a ratio of 270\,\%, $i = 3\,\%$ and $\gamma = 9\,\%$, find next year's ratio with a balanced primary budget.
\item In the same year, what \omterm{def:g12:fiscal-policy-and-debt:primary-balance}{primary balance} would have kept the ratio constant?
\item Was the interest rate above or below nominal growth on average? What does the recurrence then predict?
\end{enumerate}

\medskip
\noindent\textbf{Part II --- The decomposition.}
\begin{enumerate}[resume]
\item How many points did interest add?
\item How many points did real growth and inflation remove?
\item Find the net effect of nominal growth over the interest rate, in points.
\item Find its share of the fall.
\item What is in the row ``other''? Find its share.
\end{enumerate}

\medskip
\noindent\textbf{Part III --- Reading it.}
\begin{enumerate}[resume]
\item Inflation removed more than real growth. Why does inflation lower the ratio of a debt with fixed interest rates?
\item Who bore that part of the cost?
\item Why is the long-term rate only an approximation of the rate paid on the debt?
\item The long rate rose from 2.6\,\% in 1946 to 14.6\,\% in 1975. Why might lenders have asked for more as the period went on?
\item From the opening chart, the American ratio fell from 118.9\,\% in 1946 to 33.6\,\% in 1975. Find the fall in points.
\end{enumerate}

\medskip
\noindent\textbf{Part IV --- Lessons.}
\begin{enumerate}[resume]
\item Which of the four exits did this record use, and which not?
\item Which exit did Greece use in 2012?
\item With $i$ below $\gamma$, can a government run primary deficits for ever without the ratio exploding? Use the recurrence.
\item Why does the record not show that the same fall could be repeated with low inflation?
\item State the result of the problem: the share of the fall due to nominal growth above the interest rate.
\end{enumerate}
\end{problem}
